\chapter{Diskuze a prezentace výsledků}
Experimentální vyhodnocení ověřilo plnou životaschopnost navrženého systému na nízkoenergetické okrajové platformě (Edge AI) a díky sjednocení detekční architektury (\textbf{YOLOv8m}), prahových hodnot i polygonální oblasti zájmu (ROI) umožnilo přímou konfrontaci s~referenčním řešením na grafické kartě (GPU).

Z~hlediska celkového sčítání dopravního proudu dosáhly obě testované platformy v~náročných zimních podmínkách (husté sněžení, snížený kontrast) velmi blízkých a provozně spolehlivých výsledků. Referenční GPU řešení s~modelem v~plné přesnosti (FP32/FP16) dosáhlo celkové poměrové přesnosti sčítání \textbf{$\text{Accuracy}_{\text{count}} = 94{,}07~\%$} (111 detekovaných objektů oproti 118 skutečným, při makroskopicky odvozených hodnotách $\text{Precision} = 99{,}10~\%$, $\text{Recall} = 93{,}22~\%$ a $\text{F1} = 96{,}07~\%$). Navržené Edge řešení s~akcelerátorem Hailo-8 a 8bitově kvantizovaným modelem (INT8) dosáhlo poměrové přesnosti \textbf{$\text{Accuracy}_{\text{count}} = 92{,}37~\%$} (127 detekovaných objektů, $\text{Precision} = 89{,}76~\%$, $\text{Recall} = 96{,}61~\%$ a $\text{F1} = 93{,}06~\%$). Detailní rozbor výsledků ukázal tři klíčové technické a metodické poznatky:


\begin{itemize}
    \item \textbf{Kompromis mezi kvantizací (INT8) a plnou přesností (FP32/FP16):} Porovnání totožného modelu YOLOv8m na obou platformách prokázalo, že post-tréninková kvantizace do 8bitové celočíselné aritmetiky snižuje výslednou přesnost sčítání o~pouhých \textbf{1,70 procentního bodu}, přičemž průměrná míra spolehlivosti (\textit{confidence}) snímkových detekcí zůstává prakticky nedotčena (vážený průměr 0,7730 na Hailo-8 oproti 0,7414 na GPU). Obě platformy se však liší charakterem drobných chyb. Na akcelerátoru Hailo-8 docházelo při zákrytech vozidel za sloupem veřejného osvětlení k~mírnému kolísání rozměrů kvantizovaných ohraničujících rámečků, což vedlo k~občasné fragmentaci stopy (\textit{ghostingu}) v~modulu \texttt{hailotracker} a mírnému nadhodnocení počtu osobních automobilů (120 oproti 108, tj. $+12$), zatímco 3 menší nákladní vozy byly klasifikovány jako osobní automobily. Naopak na GPU dokázal model v~plné přesnosti bezchybně rozlišit všech 6 nákladních vozidel a díky dvoufázové asociaci algoritmu ByteTrack potlačil falešné štěpení stop ($\text{FP}=1$), avšak při přísném prahu inicializace stopy ($0{,}60$) zůstalo 7 hůře viditelných osobních vozidel v~hustém sněžení pod prahem detekce (101 oproti 108, tj. $-7$).
    
    \item \textbf{Přínos sjednocené polygonální ROI zóny:} Nasazení identického pětiúhelníkového polygonu v~knihovně Shapely na obou platformách odstranilo dřívější anomálie spojené se změnou třídy vozidla během dlouhého průjezdu širokým obdélníkem (\textit{class switching}). Polygon přesně vymezil měřený profil vozovky na výjezdu z~okružní křižovatky, což vedlo ke shodnému počtu detekovaných autobusů (4 na obou platformách oproti GT~3) i ke shodnému odfiltrování chodce pohybujícího se po chodníku mimo vozovku ($\text{FN}=1$ na obou platformách).
    
    \item \textbf{Vliv architektury datovodu na propustnost (29,39~FPS vs. 25,64~FPS) a metodické limity:} Významným zjištěním je, že okrajové zařízení Raspberry Pi~5 s~čipem Hailo-8 dosáhlo vyšší propustnosti (\textbf{29,39~FPS}, 34,02~ms/snímek) než referenční GPU pipeline (\textbf{25,64~FPS}, 39,00~ms/snímek). To potvrzuje vysokou efektivitu datově řízené architektury čipu Hailo-8 a nativního C/C++ řetězce GStreamer oproti režii zpracování obrazu a zápisu videa v~jazyce Python. Při hodnocení je zároveň nutné zohlednit dvě metodická omezení: zaprvé, prvek \texttt{queue leaky=downstream} v~GStreameru při čtení ze souboru zahodil 156 mezisnímků (4\,957 zpracovaných snímků oproti 5\,113 snímkům v~bezztrátové frontě na PC), a zadruhé, matice záměn byla vzhledem k~absenci snímkové anotace souřadnic odvozena makroskopicky z~koncových součtů tříd ($\text{TP} = \min(\text{GT}, N_{\text{sys}})$). Pro plně deterministické srovnání na úrovni jednotlivých snímků v~budoucích pracích bude vhodné vyhodnotit sledování pomocí standardizovaných metrik MOTA a IDF1 nad snímkově anotovaným datasetem.
\end{itemize}

Z~pohledu praktického nasazení v~reálné infrastruktuře chytrých měst představuje kombinace Raspberry Pi~5 a akcelerátoru Hailo-8 jednoznačně výhodnější variantu. Při zlomkové spotřebě elektrické energie ($\sim 8\text{--}12~\text{W}$ oproti $> 100~\text{W}$ u~sestavy s~dedikovanou GPU) a vyšší propustnosti 29,39~FPS dokáže autonomně a v~reálném čase zpracovávat obraz přímo u~kamery v~souladu s~principy GDPR (\textit{Privacy by Design}), přičemž přesnost sčítání přesahující 92~\% v~náročných zimních podmínkách zaostává za plným GPU řešením o~méně než dvě procenta a je pro dlouhodobé sledování intenzity dopravy plně dostačující.